\section{Structure coefficients at every positive length}
\label{sec:asym_unblocked_coefficients}

Structure coefficients in an expansion of an arbitrary tensor over
inequivalent normal tensors are power sums of fixed finite multisets of
complex weights, at every positive length and without blocking. Normality
means injectivity after some blocking (Definition~\ref{def:normal}); the
tensors need not be normalized. Separation excludes gauge equivalence up to
any nonzero complex scalar, not only a unit-modulus phase.

\begin{theorem}[Uniqueness of nonzero weights from tail power sums]
    \label{thm:asym_tail_weights}
    \lean{Multiset.eq_of_notMem_zero_of_forall_sum_map_pow_eq,
        Multiset.eq_zero_of_notMem_zero_of_forall_sum_map_pow_eq_zero}
    \leanok
    Let $\Lambda$ and $\Lambda'$ be finite multisets in $\C\setminus\{0\}$.
    If their power sums agree for all sufficiently large integers $N$, then
    $\Lambda=\Lambda'$, including multiplicities. In particular, if the
    power sums of $\Lambda$ vanish for all sufficiently large $N$, then
    $\Lambda$ is empty.
\end{theorem}
\begin{proof}
    \leanok
    List the distinct entries of the two multisets as
    $z_0,\ldots,z_{r-1}$, and let $b_j$ be the difference of their
    multiplicities. For a threshold $N_0$ and each $0\leq i<r$,
    \begin{align}
        \sum_{j=0}^{r-1}(b_jz_j^{N_0})z_j^i & = 0.
    \end{align}
    The Vandermonde matrix is invertible, so $b_jz_j^{N_0}=0$.
    Since $z_j\neq0$, all multiplicities agree. Comparison with the empty
    multiset proves the vanishing assertion.
\end{proof}

\begin{theorem}[Power-sum coefficients without blocking]
    \label{thm:asym_unblocked_coefficients}
    \lean{MPSTensor.exists_unblocked_powerSum_coeff_unique,
        MPSTensor.exists_unblocked_powerSum_coeff}
    \leanok
    \uses{def:normal}
    Let $B$ be an MPS tensor of bond dimension $D_B$, and let
    $\{M_\gamma\}_{\gamma\in\Gamma}$ be a finite family of normal MPS
    tensors on the same physical alphabet, with $D_\gamma>0$.
    For distinct $\gamma,\delta$ of equal bond dimension, assume there
    are no $X\in\GL(D_\gamma,\C)$ and $z\in\C\setminus\{0\}$ such that
    $M_\delta^i=zXM_\gamma^iX^{-1}$ for every physical index $i$.
    Suppose that for every $N>0$ there are coefficients $c_\gamma^{(N)}$
    satisfying
    \begin{align}
        \ket{V^{(N)}(B)}
         & =\sum_{\gamma\in\Gamma}c_\gamma^{(N)}
        \ket{V^{(N)}(M_\gamma)}.
        \label{eq:asym_unblocked_span}
    \end{align}
    Then there are integers $n_\gamma\geq0$ and nonzero complex weights
    $\mu_{\gamma k}$, indexed by $0\leq k<n_\gamma$, such that for every
    $N>0$,
    \begin{align}
        \ket{V^{(N)}(B)}
         & =\sum_{\gamma\in\Gamma}
        \left(\sum_{k=0}^{n_\gamma-1}\mu_{\gamma k}^{N}\right)
        \ket{V^{(N)}(M_\gamma)}.
        \label{eq:asym_unblocked_expansion}
    \end{align}
    The multiplicities satisfy $\sum_{\gamma\in\Gamma}n_\gamma D_\gamma\leq D_B$.
    There is a positive threshold $N_0$ such that, for every $N\geq N_0$,
    any coefficients satisfying~(\ref{eq:asym_unblocked_span}) obey
    $c_\gamma^{(N)}=\sum_{k=0}^{n_\gamma-1}\mu_{\gamma k}^{N}$.
    The weight multisets are unique: any finite families of nonzero
    weights giving an expansion of the form~(\ref{eq:asym_unblocked_expansion})
    at all sufficiently large lengths have, for each $\gamma$, exactly
    the same entries with the same multiplicities. Empty weight
    multisets are allowed.
\end{theorem}

\begin{center}
    % Formula: the component of (\ref{eq:asym_unblocked_expansion}) at
    % the configuration (i_1,\ldots,i_N):
    % \tr(B^{i_1}\cdots B^{i_N})=\sum_\gamma(\sum_k\mu_{\gamma k}^N)
    % \tr(M_\gamma^{i_1}\cdots M_\gamma^{i_N}).
    % Source: Theorem~\ref{thm:asym_unblocked_coefficients}.
    % Ink-to-index map: the left ring is the periodic chain of $B$, the
    % right ring the periodic chain of $M_\gamma$; the site in column $n$
    % carries the physical letter $i_n$.
    % Contraction set: every horizontal virtual bond, closed by the
    % periodic trace; the physical legs remain open. The sum over
    % $\gamma$ and the power-sum coefficient are scalars, not
    % contraction edges.
    % Boundary signature: (virtual-west=0 | virtual-east=0 |
    % physical-up=3 | physical-down=0) for the schematic on each side.
    \tenkzkernel
    \begin{tenkz}[rows={ket}, cols=4, physical=up, boundary=periodic]
        \tn[label pos=s, ports={90:physical:$i_1$}]{B} &
        \tn[label pos=s, ports={90:physical:$i_2$}]{B} &
        \tn[skin=dots, ports={180:virtual, 0:virtual}]{} &
        \tn[label pos=s, ports={90:physical:$i_N$}]{B}
    \end{tenkz}
    \quad $=\;\displaystyle\sum_{\gamma\in\Gamma}
        \Bigl(\sum_{k=0}^{n_\gamma-1}\mu_{\gamma k}^{N}\Bigr)$ \quad
    \begin{tenkz}[rows={ket}, cols=4, physical=up, boundary=periodic]
        \tn[label pos=s, ports={90:physical:$i_1$}]{M_\gamma} &
        \tn[label pos=s, ports={90:physical:$i_2$}]{M_\gamma} &
        \tn[skin=dots, ports={180:virtual, 0:virtual}]{} &
        \tn[label pos=s, ports={90:physical:$i_N$}]{M_\gamma}
    \end{tenkz}
\end{center}
\begin{proof}
    \leanok
    \uses{thm:irred_decomp,thm:periodic_overlap_dichotomy,
        thm:periodic_basis_eventually_li,thm:asym_tail_weights}
    Apply Theorem~\ref{thm:irred_decomp}, normalize each nonzero irreducible
    factor to a periodic tensor, and group factors whose periodic vectors
    are proportional by a fixed nonzero scalar raised to the length.
    The irreducible-form framework is described in
    Ref.~\cite{DeLasCuevas2017Irreducible}. This gives separated periodic
    representatives $A_j$ of dimensions $d_j$, positive multiplicities
    $r_j$, and nonzero weights $a_{jq}$ such that, for every $N>0$,
    \begin{align}
        \ket{V^{(N)}(B)}
                      & =\sum_j\left(\sum_{q=0}^{r_j-1}a_{jq}^{N}\right)
        \ket{V^{(N)}(A_j)}, \label{eq:asym_periodic_expansion}           \\
        \sum_j r_jd_j & \leq D_B.
    \end{align}
    Grouping preserves the dimension count. Indeed, if periodic tensors
    $A,A'$ of periods $m,m'$ have such proportional vectors, then
    $|O_{AA'}(N)|^2=O_{AA}(N)O_{A'A'}(N)$.
    Different dimensions would force the left side to tend to zero by
    Theorem~\ref{thm:periodic_overlap_dichotomy}, whereas along common
    multiples of the periods the right side tends to $mm'>0$.

    Normalize the targets to period-one tensors $C_\gamma$ and combine
    them with the representatives not related to any target by an invertible gauge
    and a nonzero scalar.
    The combined family has no repeated blocks. On sufficiently large
    multiples $pn$ of a common period, it is linearly independent by
    Theorem~\ref{thm:periodic_basis_eventually_li}.
    Subtracting~(\ref{eq:asym_unblocked_span}) from
    (\ref{eq:asym_periodic_expansion}), after absorbing matched terms
    into the targets, forces every unmatched representative to satisfy
    $\sum_{q=0}^{r_j-1}(a_{jq}^{p})^n=0$ for all sufficiently large $n$.
    Theorem~\ref{thm:asym_tail_weights} contradicts $r_j>0$.
    Thus every representative matches a target.

    Choose its matching label $\gamma(j)$ and a nonzero $\theta_j$ so that
    $\ket{V^{(N)}(M_{\gamma(j)})}
        =\theta_j^N\ket{V^{(N)}(A_j)}$.
    For each target, collect the weights $a_{jq}/\theta_j$ over its
    matching representatives. This proves~(\ref{eq:asym_unblocked_expansion}).
    Since $d_j=D_{\gamma(j)}$, the dimension bound follows.
    Normalize the normal targets and use their eventual linear independence
    to compare arbitrary expansion coefficients with these power sums.
    The nonzero normalization scalars can be cancelled, giving a positive
    threshold for coefficient uniqueness. Theorem~\ref{thm:asym_tail_weights}, applied to each
    target, then proves uniqueness against any tail expansion.
\end{proof}

\begin{theorem}[Power-sum structure coefficients for operator closure]
    \label{thm:asym_operator_coefficients}
    \lean{MPOTensor.exists_operatorProduct_powerSum_coeff_unique}
    \leanok
    \uses{def:normal}
    Let $\{M_\gamma\}_{\gamma\in\Gamma}$ be a finite family of MPO
    tensors of common physical dimension $d$, with $D_\gamma>0$. Regard the pair of physical indices
    $(i,j)$ as one MPS index. Assume these MPS tensors are normal and,
    for distinct labels of equal bond dimension, are not related by
    $M_\delta^{ij}=zXM_\gamma^{ij}X^{-1}$ for all $i,j$, with $z\neq0$
    and $X$ invertible. Fix $\alpha,\beta\in\Gamma$ and write $O_L(M)$
    for the periodic operator generated by $M$. Suppose
    $O_L(M_\alpha)O_L(M_\beta)$ belongs to
    $\spn_{\C}\{O_L(M_\gamma):\gamma\in\Gamma\}$ for every $L>0$.

    There are complex diagonal matrices $\chi_{\alpha,\beta,\gamma}$
    of sizes $n_\gamma\geq0$, with nonzero diagonal entries, such that
    for every $L>0$,
    \begin{align}
        O_L(M_\alpha)O_L(M_\beta)
         & =\sum_{\gamma\in\Gamma}
        \tr(\chi_{\alpha,\beta,\gamma}^{L})O_L(M_\gamma).
        \label{eq:asym_operator_expansion}
    \end{align}
    Their sizes satisfy $\sum_{\gamma\in\Gamma}n_\gamma D_\gamma
        \leq D_\alpha D_\beta$.
    Beyond a positive threshold, any coefficients of this operator
    expansion equal $c_{\alpha,\beta,\gamma}^{(L)}
        =\tr(\chi_{\alpha,\beta,\gamma}^{L})$.
    Any other finite families of nonzero diagonal entries giving such an
    expansion at all sufficiently large lengths agree target by target,
    including multiplicities.

    This uses the structure-coefficient notation of
    Ref.~\cite[Section~4.5]{Cirac2016MPDO_arXiv}.
\end{theorem}

\begin{center}
    % Formula: (\ref{eq:asym_operator_expansion}),
    % O_L(M_\alpha)O_L(M_\beta)
    % =\sum_\gamma\tr(\chi_{\alpha,\beta,\gamma}^L)O_L(M_\gamma).
    % Source: Theorem~\ref{thm:asym_operator_coefficients}; the closed
    % fusion identity of \cite[Section~4.5]{Cirac2016MPDO_arXiv}, drawn
    % as in Figure~\ref{fig:mpdo_fusion_trace_power}.
    % Ink-to-index map: on the left, the upper row is $M_\alpha$ with
    % output legs $i_1,\ldots,i_L$ and the lower row is $M_\beta$ with
    % input legs $j_1,\ldots,j_L$; the vertical pairleg at each site is
    % the summed middle physical index. On the right, the upper loop of
    % boxes labelled $\chi$ is $\tr(\chi_{\alpha,\beta,\gamma}^{L})$ and
    % the lower loop is $O_L(M_\gamma)$.
    % Contraction set: on the left, the pairlegs at every displayed site
    % and the horizontal virtual bonds of each row, both rows closed by
    % the periodic trace; on the right, the virtual bonds of each loop.
    % The two loops share no index, so their juxtaposition is the
    % product of a scalar and an operator; the sum over $\gamma$ is a
    % sum of pictures, not a contraction edge. The left panel authors
    % its bonds and leaves the ellipsis column bare, since the frame
    % would otherwise mint a pairleg stub between the two ellipses.
    % Boundary signature: (virtual-west=0 | virtual-east=0 |
    % physical-up=3 | physical-down=3) for the schematic on each side.
    \tenkzkernel
    \begin{tenkz}[rows={op,op}, cols=4, west=trace, east=trace, bonds=none]
        \tn[skin=mpo, ports={90:physical:$i_1$, 270:physical}]{M_\alpha} &
        \tn[skin=mpo, ports={90:physical:$i_2$, 270:physical}]{M_\alpha} &
        \tn[skin=dots]{} &
        \tn[skin=mpo, ports={90:physical:$i_L$, 270:physical}]{M_\alpha} \\
        \tn[skin=mpo, ports={90:physical, 270:physical:$j_1$}]{M_\beta} &
        \tn[skin=mpo, ports={90:physical, 270:physical:$j_2$}]{M_\beta} &
        \tn[skin=dots]{} &
        \tn[skin=mpo, ports={90:physical, 270:physical:$j_L$}]{M_\beta}
        \tnwire{(1,1)}{(1,2)} \tnwire{(1,2)}{(1,3)} \tnwire{(1,3)}{(1,4)}
        \tnwire{(2,1)}{(2,2)} \tnwire{(2,2)}{(2,3)} \tnwire{(2,3)}{(2,4)}
        \tnwire{(1,1)}{(2,1)} \tnwire{(1,2)}{(2,2)} \tnwire{(1,4)}{(2,4)}
    \end{tenkz}
    \quad $=\;\displaystyle\sum_{\gamma\in\Gamma}$ \quad
    \begin{tabular}[c]{@{}c@{}}
        {\begin{tenkz}[cols=4, west=trace, east=trace]
             \tn[skin=box]{\chi} & \tn[skin=box]{\chi} & \tn[skin=dots]{} &
             \tn[skin=box]{\chi}
         \end{tenkz}} \\
        {\begin{tenkz}[cols=4, west=trace, east=trace]
             \tn[skin=mpo,
                 ports={90:physical:$i_1$, 270:physical:$j_1$}]{M_\gamma} &
             \tn[skin=mpo,
                 ports={90:physical:$i_2$, 270:physical:$j_2$}]{M_\gamma} &
             \tn[skin=dots]{} &
             \tn[skin=mpo,
                 ports={90:physical:$i_L$, 270:physical:$j_L$}]{M_\gamma}
         \end{tenkz}}
    \end{tabular}
\end{center}
\begin{proof}
    \leanok
    \uses{thm:mpdo_operator_product_mpo,thm:asym_unblocked_coefficients}
    Form the product tensor $P^{ij}=\sum_s M_\alpha^{is}\otimes M_\beta^{sj}$,
    of bond dimension $D_\alpha D_\beta$.
    Theorem~\ref{thm:mpdo_operator_product_mpo} gives
    $O_L(P)=O_L(M_\alpha)O_L(M_\beta)$.
    Equality of operators is equivalent, entry by entry, to equality of
    MPS coefficients on the combined physical alphabet.
    Apply Theorem~\ref{thm:asym_unblocked_coefficients} to $P$ and take
    $\chi_{\alpha,\beta,\gamma}$ to have diagonal entries
    $\{\mu_{\gamma k}\}_{k=0}^{n_\gamma-1}$.
    The same entrywise equivalence transfers the expansion, eventual
    coefficient uniqueness, dimension bound, and uniqueness of nonzero
    entries to operators.
\end{proof}
